% چکیده فارسی
\chapter*{چکیده}
\phantomsection
\addcontentsline{toc}{chapter}{چکیده}

حل عددی معادلات دیفرانسیل جزئی غیرخطی یکی از مسائل بنیادی و پرکاربرد در محاسبات علمی و مهندسی به شمار می‌رود. روش‌های عددی کلاسیک، با وجود دقت و مبانی نظری استوار، به گسسته‌سازی و تولید شبکه محاسباتی وابسته‌اند و در مواجهه با ابعاد بالا یا هندسه‌های پیچیده با چالش‌های جدی روبه‌رو می‌شوند. از سوی دیگر، مدل‌های یادگیری عمیقِ صرفاً داده‌محور، افزون بر نیاز به حجم بالایی از داده‌های آموزشی، تضمینی برای ارضای قوانین و اصول بقای فیزیکی ارائه نمی‌دهند. چارچوب «شبکه‌های عصبی آگاه از فیزیک» (\lr{PINN}) به‌عنوان رهیافتی نوآورانه برای پیوند دادن دانش پیشین فیزیکی (معادلات دیفرانسیل حاکم بر سیستم) با توانمندی تقریب‌زدن شبکه‌های عصبی عمیق مطرح گردیده است. هدف این پایان‌نامه، مطالعه نظری، پیاده‌سازی محاسباتی و ارزیابی جامع این چارچوب در حل مسائل مستقیم و معکوس معادلات دیفرانسیل جزئی غیرخطی است.

در این پژوهش، پس از تبیین مبانی ریاضی معادلات دیفرانسیل جزئی، روش‌های عددی کلاسیک، یادگیری عمیق و سازوکار مشتق‌گیری خودکار، فرمول‌بندی شبکه‌های آگاه از فیزیک در دو گونه مدل زمان‌پیوسته و زمان‌گسسته تشریح شده است. در بخش عملی، مدل زمان‌پیوسته برای معادله غیرخطی برگرز با استفاده از زبان پایتون و کتابخانه جکس پیاده‌سازی گردید. همچنین یک حل‌گر مرجع مستقل مبتنی بر روش طیفی فوریه با گام‌زنی زمانی مرتبه چهار نمایی (\lr{ETDRK4}) برای تولید جواب دقیق توسعه یافت و صحت عملکرد آن با جواب نیمه‌تحلیلی تبدیل کول-هوپف راستی‌آزمایی شد.

در مسئله مستقیم، میدان جواب در کل دامنه زمانی-مکانی تنها با بهره‌گیری از ۲۰۰ نقطه داده روی شرایط اولیه و مرزی با خطای نسبی \FwdRelErr\ در نرم $L_2$ نسبت به حل مرجع بازسازی شد؛ تحلیل مکانی خطا نشان داد که بیشینه عدم‌قطعیت در لایه باریک جبهه شوک متمرکز است. در مسئله معکوس (کشف پارامترهای نامعلوم معادله)، ضرایب انتقال و پخش از روی ۵۰۰۰ مشاهده پراکنده در دو سناریوی داده تمیز و داده آلوده به یک درصد نویز گاوسی شناسایی شدند: در سناریوی بدون نویز ضریب انتقال با خطای نسبی \LamIErrClean\ و ضریب پخش با خطای \LamIIErrClean، و در سناریوی نویزی به‌ترتیب با خطای \LamIErrNoisy\ و \LamIIErrNoisy\ برآورد گردیدند. یافته‌های این پژوهش مؤید کارایی داده‌ای چشمگیر چارچوب شبکه‌های آگاه از فیزیک و پایداری رضایت‌بخش آن در برابر اغتشاش داده‌هاست و نشان می‌دهد که چالش‌های اصلی روش به تفکیک لایه‌های با گرادیان تند و حساسیت ذاتی در شناسایی ضرایب کوچک پخش بازمی‌گردد.

\textbf{کلیدواژه‌ها:} شبکه‌های عصبی آگاه از فیزیک، معادلات دیفرانسیل جزئی غیرخطی، مسئله مستقیم، مسئله معکوس، معادله برگرز، مشتق‌گیری خودکار، روش طیفی فوریه.
